%-----------------------Homework------------------------------------
%-------------------Arman Shokrollahi---------------------------------
%---------------------Coding Theory-------------------------------

\documentclass[a4 paper]{article}
% Set target color model to RGB
\input{preamble.tex}

\begin{document}
\homework{Práctia \#4}{Abril. 22 al 23 de 2019}{MSc. Fausto M. Lagos S.}{}{Aquí va su nombre}{}

En esta práctica se desarrollarán tres planteamientos cada uno con una valoración de dos puntos. El objetivo es afianzar sus conocimientos antes de la Prueba de Evaluación Continua correspondiente, puede utilizar todo el material bibliográfica a su disposición y también preguntar todo lo que considere necesario. Preste atención al margen derecho donde encontrará premios y bonificaciones adicionales.

\vspace{5mm}

\problem{1} \textbf{2 puntos} - De El Cálculo de Louis Leithold 7ed, pág 34 lea detenidamente el Ejemplo 4 y resuelva los enunciados 34 y 36 de la página 38.
\vskip 3mm
\solution{} 

\vskip 3mm

\problem{2} \textbf{2 puntos} - Del Cálculo con trascendentes temprana 4ed de Dennis G. Zill, lea el Ejemplo 7 de la página 71 y resuelva los siguientes enunciados:
\begin{itemize}
    \item[(a)] \[\lim_{x\to 0} \frac{x}{\sin (3x)}\]
    \item[(b)] \[\lim_{x\to 0} \frac{\tan x}{x}\]
\end{itemize}
\vskip 3mm
\solution{}
\vspace{30mm}

\problem{3} \textbf{2 puntos} - \marginnote{\includegraphics[scale = 1]{overleaf}} Dibuje la gráfica de alguna función $f$ que satisfaga las siguientes condiciones: $f(-2)\neq f(2)$; $\displaystyle{\lim_{x\to -2}f(x) \neq f(-2)}$; $\displaystyle{\lim_{x\to 2}f(x) \neq f(2)}$; $\displaystyle{\lim_{x\to a}f(x) = f(a)}$ si $a\neq \pm 2$; el contradominio de $f$ es el intervalo cerrado $[-3, 3]$.
\vskip 3mm
\solution{}
\vspace{30mm}

\end{document}
